\documentclass[11pt]{article}
\usepackage[T1]{fontenc}
\usepackage{lmodern,microtype,amsmath,amssymb,amsthm,booktabs,array}
\usepackage[letterpaper,margin=0.95in,top=0.88in,bottom=0.9in]{geometry}
\usepackage[dvipsnames]{xcolor}
\usepackage[colorlinks=true,linkcolor=MidnightBlue,urlcolor=MidnightBlue]{hyperref}
\usepackage{fancyhdr}
\newtheorem{theorem}{Theorem}[section]
\newtheorem{lemma}[theorem]{Lemma}
\newtheorem{proposition}[theorem]{Proposition}
\newtheorem{corollary}[theorem]{Corollary}
\theoremstyle{remark}\newtheorem{remark}[theorem]{Remark}
\setlength{\headheight}{14pt}
\pagestyle{fancy}\fancyhf{}\fancyhead[L]{\small\textsc{Applications of Dirichlet GRH}}\fancyhead[R]{\small Research note}\fancyfoot[C]{\thepage}
\urlstyle{same}\raggedbottom\setlength{\emergencystretch}{2em}
\hypersetup{pdftitle={Sparse CM equations, prime-generating constants, and additive representations}}
\begin{document}
\begin{center}
{\small\scshape Number theory\enspace /\enspace Research note\par}
\vspace{.65em}
{\LARGE\bfseries Sparse CM equations, prime-generating constants,\par and additive representations\par}
\vspace{.45em}
{\large Applications of the Dirichlet GRH theorem in \texttt{grh-lean}}\\[.55em]
25 November 2026
\end{center}
\begin{abstract}
For a singular modulus of order discriminant $\Delta$, let $s(\Delta)$ be the fewest terms in any nonzero rational polynomial that vanishes at it. The verified Dirichlet GRH theorem supplies the hypothesis of the published proof of Riffaut's trinomial conjecture. Combining quadratic prime splitting with separated absolute values of conjugates gives the stronger asymptotic bound $s(\Delta)\gg |\Delta|^{1/4}/\sqrt{\log |\Delta|}$. An analytic cutoff and exhaustive finite checks exclude four-term equations in every degree at least six and reduce degrees four and five to 79 explicit candidates. One unconditional integer certificate requires at least 7,959 terms for $\Delta=1-2^{61}$. In a second application we correct the least prime-generating constant claimed for exponent $9/4$ in an identifiable preprint example and certify its first 42 decimal places. The same construction gives rigorous digits for four exponents down to $2.01$. A third application classifies all exceptions to a prime plus a positive multiple of a prime cube or fourth power. Finally, a variant of the repository's internal analytic argument detects zeros from upper bounds sampled on thin sets of twists and scales. The analytic deductions are paper proofs. Integer certificates accompany the finite results.
\end{abstract}
\section{The theorem being used}
We inspect the readable edition of \href{https://github.com/lsernau/grh-lean/tree/72fead102f663d5c7a407964a13dc0cf9d977207}{\texttt{lsernau/grh-lean} at \texttt{72fead10}}. Its \href{https://github.com/lsernau/grh-lean/blob/72fead102f663d5c7a407964a13dc0cf9d977207/grh/src/Statements.lean}{\texttt{Statements.lean}} imports Mathlib alone. The declarations in \href{https://github.com/lsernau/grh-lean/blob/72fead102f663d5c7a407964a13dc0cf9d977207/grh/src/Main.lean}{\texttt{Main.lean}} prove Mathlib's ordinary Riemann hypothesis, give its explicit zero statement, and establish the following three Dirichlet assertions. For every modulus $N\ge1$ and every Dirichlet character $\chi$ of that modulus,
\[
 \Re s>1/2,\quad (\chi\ne\mathbf1\ \text{or}\ s\ne1)
        \quad\Longrightarrow\quad L(s,\chi)\ne0.
\]
Every zero in $0<\Re s<1$ is on the critical line. The third declaration specializes the line statement to primitive characters. This scope includes quadratic characters of all fundamental discriminants and the real and complex characters needed to count primes in arbitrary reduced residue classes. Ordinary RH alone suffices for the prime-generating constants.

The readable edition pins Lean 4.35.0 and Mathlib \texttt{56a00585}. Andrew Sutherland's verification record reports sandboxed builds, kernel replays of every module, and independent checks of the exported theorems in both editions. The reported axioms are \texttt{propext}, \texttt{Classical.choice} and \texttt{Quot.sound}. The separate challenge file supplies statements for comparison. The upstream README credits Levent Alp\"oge's proof with Claude as the basis for the shorter development. For this note we checked the statements and 15 source-blob identities. We rely on the recorded independent kernel checks and did not rebuild the full project. The results below have not been added to its Lean proof. See the \href{https://math.mit.edu/~drew/shortgrh/}{independent verification record} and \href{https://github.com/lsernau/grh-lean/blob/72fead102f663d5c7a407964a13dc0cf9d977207/grh/README.md}{readable edition README}.

\paragraph{How the proof reaches a zero-free half-plane.}
The formal proof introduces an auxiliary quadratic twist and a nonnegative weighted prime sum. It bounds the same lattice of automorphic forms in two ways. Short integral forms and a second moment give enough independent vectors for an upper bound. A twisted determinant gives a lower bound on their corrected covolume. Canceling the leading terms leaves an upper bound on the remainder of the prime sum. A zero-density result supplies an auxiliary twist free of zeros in a chosen rectangle. A one-sided Landau argument then excludes an untwisted zero there. The residue and interpolation steps, which can be used while weakening the sampling hypotheses, are examined in Section~\ref{sec:sparse}. The theta lifts, lattices and choice of twists used to establish the original prime-sum bound depend on the particular arithmetic family. See the \href{https://github.com/lsernau/grh-lean/blob/72fead102f663d5c7a407964a13dc0cf9d977207/grh/GUIDE.md}{upstream reading guide}.

\paragraph{Relation to the three incoming reports.}
We compared three archives in \href{https://github.com/VladimirReshetnikov/ProveIt/tree/1e1728ad221d69e48fea09bf754b73422af44943/docs/incoming}{the incoming directory at \texttt{1e1728ad}}.
\begin{center}\small
\begin{tabular}{@{}p{.30\linewidth}p{.65\linewidth}@{}}
\toprule
Archive & Subjects already developed\\ \midrule
\texttt{grh-lean-transfers.zip} & Quintic coefficient fields of CM newforms, primitive roots and norm-Euclidean fields, and a cubic prime-sum test.\\
\texttt{grh-transfer.zip} & An abelian field census, idoneal numbers and applications, and the scope of Kummer-field reductions.\\
\texttt{grh\_transfers.zip} & Class-group exponent restrictions, a mod-four weighted prime race, and a discriminant $-23$ cubic test.\\
\bottomrule
\end{tabular}\end{center}
The CM objects here are the singular values of $j$ themselves. This is a different invariant from the fields generated by Fourier coefficients of CM newforms. The growth bound, finite classifications and minimum-constant certificates below are deductions of this analysis. We claim no bibliographic priority. Named conditional implications retain their original attribution.

\section{Riffaut's conjecture and sparse equations for singular moduli}

A singular modulus is a value $\alpha=j(\tau)$ of the modular invariant at
an imaginary quadratic argument in the upper half-plane. The corresponding
elliptic curve has complex multiplication by an order of discriminant
$\Delta<0$. The conjugates of $\alpha$ over $\mathbb Q$ are exactly the
singular moduli of discriminant $\Delta$, and their number is the order's
class number $h(\Delta)$. A polynomial's support size is its number of
nonzero coefficients. Write $s(\Delta)$ for the minimum support size of a
nonzero rational polynomial vanishing at one, and hence every, such modulus.

\begin{theorem}[Transfer to Riffaut's conjecture]
Accepting the Dirichlet GRH theorem exported by the repository,
$s(\Delta)\ge4$ whenever $h(\Delta)\ge3$.
\end{theorem}
\begin{proof}
Suppose $\alpha$ satisfies a relation of support at most three.
It is nonzero since it has degree at least three. Dividing by its smallest
power of $\alpha$ and its leading coefficient gives
$\alpha^u+A\alpha^v+B=0$, with $u>v>0$, $A,B\in\mathbb Q$ and
$B\ne0$. A binomial relation of degree at least three permits $A=0$.

Bilu, Luca and Pizarro-Madariaga call such a $\Delta$ a trinomial
discriminant. Their Theorems 1.3 and 1.4 imply
$m=|\Delta|>10^{11}$ and $\Delta=-p$ or $-pq$ for distinct odd primes
in the second case, with $\Delta$ fundamental. Hence the Kronecker character
$\chi=(\Delta/\cdot)$ is primitive, real, and odd, of conductor $m$ with at
most two distinct prime factors. Their Corollary 5.2 forces every prime $p$
with $\chi(p)=1$ to satisfy $p>4\sqrt m/\log m$. Their Proposition 9.2,
under Dirichlet GRH, produces a prime with $\chi(p)=1$ and
$p\le4\sqrt m/\log m$ for this very character, since $m\ge10^{10}$.
This contradiction is their Theorem 9.1. Its zero-location input is supplied
by the upstream theorems for the Riemann zeta-function and primitive Dirichlet
characters. The results about singular moduli and the explicit constants
are the published authors' work.
\end{proof}

The two arguments meet at a specific arithmetic object. Section 9.1 of
the published proof uses the smoothed sum
$T(x,\chi)=\sum_{n\le x}\chi(n)\Lambda(n)(1/n-1/x)$.
The repository's prime remainder uses the same kernel with weight
$1-\chi_K(n)$ in the progression $1\bmod q$. Its eventual estimate has
parameters depending on the modulus. The explicit conductor threshold in
Proposition 9.2 enters through the published argument after applying GRH.
Full details appear in Y.~Bilu, F.~Luca and A.~Pizarro-Madariaga,
\emph{Trinomials, singular moduli and Riffaut's conjecture},
Ann. Sc. Norm. Super. Pisa Cl. Sci. (5) 23 (2022), 2003--2048,
\href{https://arxiv.org/pdf/2003.06547}{arXiv:2003.06547v2},
especially \S\S2, 5--7, 9.
A literature search located this earlier conditional proof. We make no claim
to be the first to observe the consequence of the new upstream theorem.

\begin{corollary}[Coefficients in the CM field]
Let $K=\mathbb Q(\sqrt\Delta)$. The minimum support size stays exactly
$s(\Delta)$ when the coefficients range over $K$.
\end{corollary}
\begin{proof}
The ring class field theorem gives
$[K(\alpha):K]=h(\Delta)=[\mathbb Q(\alpha):\mathbb Q]$.
Consequently the natural map
$K\otimes_{\mathbb Q}\mathbb Q(\alpha)\to K(\alpha)$ is an isomorphism.
For any finite choice of exponents, independence of the corresponding powers
over $\mathbb Q$ therefore persists over $K$. Conversely a rational relation
is also a relation over $K$. The same argument works for any coefficient
field linearly disjoint from $\mathbb Q(\alpha)$ over $\mathbb Q$.
\end{proof}

\subsection{A density-one extension to any fixed number of terms}

For four or more terms, the preceding published trinomial theorem leaves the
minimum support undetermined. The same geometry of CM conjugates yields a bound
which can exceed any prescribed support size. The following calculation
uses the reduced forms of Bilu--Luca--Pizarro-Madariaga, \S2 and the proof
of Proposition 3.1(6), and the root-counting argument of their Lemma 12.3.
It holds independently of the repository's GRH theorem.

\begin{proposition}[Support from coprime factorizations]
Let $\omega(n)$ be the number of distinct prime divisors of $n$. Then
\[
\begin{aligned}
 s(-4n)&\ge 1+2^{\omega(n)-1} &&(n>1),\\
 s(-N)&\ge 1+\left\lceil 2^{\omega(N)-2}\right\rceil
 &&(N>3,\ N\equiv3\pmod4).
\end{aligned}
\]
In particular every order discriminant $-4n$ with at least three distinct
primes dividing $n$ excludes all four-term rational polynomials.
\end{proposition}
\begin{proof}
There are $2^{\omega(n)-1}$ unordered coprime factorizations $n=uv$,
$u<v$. Assign each full prime power to one of two factors and identify the
two orderings. The forms $[u,0,v]$ are distinct primitive reduced forms
of discriminant $-4n$. Their CM points are $i\sqrt{v/u}$. All give positive
conjugates, since for $z$ on the positive imaginary axis the formula
$j(z)=E_4(z)^3/(q\prod_{m\ge1}(1-q^m)^{24})$, $q=e^{2\pi iz}$, has
positive numerator and denominator by the Eisenstein series expansion.
Descartes' rule gives at most $k-1$ positive roots for a polynomial with
$k$ terms. This proves the even bound.

For $N=uv$ odd, again use coprime pairs $u<v$. They are opposite modulo
$4$. If $v>3u$, take $[u,u,(u+v)/4]$. If $v<3u$, take
$[(u+v)/4,(v-u)/2,(u+v)/4]$. Equality forces $(u,v)=(1,3)$, excluded.
Each triple is integral, reduced and has discriminant $-N$. For
primitivity, an odd common prime divisor would divide both $u$ and $v$,
and the middle coefficient is odd. Each pair can be recovered from its
triple, and the two types have disjoint shapes. The associated CM points
$z=(b+i\sqrt N)/(2a)$ have either $\Re z=1/2$ or $|z|=1$. In the first
case $-\bar z=z-1$ and in the second $-\bar z=-1/z$, so modular invariance
and the real Fourier coefficients of $j$ show that $j(z)$ is real. All
these conjugates are nonzero because the unique zero of $j$ on the
modular quotient has discriminant $-3$. Of $2^{\omega(N)-1}$ real
conjugates, at least half have the same sign. Apply Descartes' rule to
$P(X)$ or $P(-X)$ to finish.
\end{proof}

\begin{corollary}[Density for order discriminants]
Fix $k\ge1$. Among negative order discriminants, counted once each in
order of absolute value, the proportion with $s(\Delta)\le k$ tends to zero.
The same holds for equations with coefficients in the corresponding CM field.
\end{corollary}
\begin{proof}
For either branch of the proposition, bounded $s$ forces bounded $\omega$.
For clarity, fix $L$ and select the odd primes up to $y$. For
integers $m$ with natural-density weights, divisibility by these primes is a
family of independent indicators, by the Chinese remainder theorem.
Independence and the individual probabilities $1/p$ persist on
$m\equiv3\pmod4$. Their sum has mean $\mu_y=\sum_{2<p\le y}1/p$
and variance at most $\mu_y$. For $\mu_y>L$, Chebyshev's inequality bounds
the upper density of $\omega(m)\le L$ by $\mu_y/(\mu_y-L)^2$ in either
population. Divergence of the prime reciprocal sum lets $y$ tend to
infinity. Apply this to $m=n$ in the even case and to $m=N$ in the odd case.
The two cases exhaust the order discriminants up to finitely many exceptions.
\end{proof}

Two examples attain the bound. At $\Delta=-120$ the complete list of
reduced primitive forms is $[u,0,30/u]$ for $u=1,2,3,5$. At $\Delta=-840$
it is $[u,0,210/u]$ for $u=1,2,3,5,6,7,10,14$. The class polynomials have
degrees $4$ and $8$, so their support sizes are at most $5$ and $9$.
The even bound proves equality for the minimum over all polynomial
multiples, of arbitrary degree. The accompanying
\texttt{verify\_cm\_support.py} checks the reduced-form lists with integer
arithmetic. The density statement weights each discriminant once. Individual
odd prime discriminants can have arbitrarily large absolute value while
giving only one coprime factor pair, so the factor-pair bound alone gives no growing lower bound in that case.

\subsection{A uniform support bound from primes splitting in the CM field}

The zero-location theorem also gives a bound that grows with the discriminant
even when the discriminant has only one prime factor. We derive it by finding
many conjugates of different absolute sizes. The exponents and coefficients
of a potential sparse equation may be arbitrarily large.

\begin{lemma}[Separated conjugate sizes]
Let $z_1,\ldots,z_t$ be nonzero complex numbers whose moduli $r_i=|z_i|$
satisfy $r_{i+1}\ge16r_i$. For any integers $0\le e_1<\cdots<e_t$,
the matrix $(z_i^{e_j})_{i,j=1}^t$ is invertible. In particular a polynomial
vanishing on all these points has more than $t$ terms unless it is zero.
\end{lemma}
\begin{proof}
Multiply column $j$ by a positive number $c_j$, starting with $c_1=1$ and
choosing $c_{j+1}/c_j=(r_jr_{j+1})^{-(e_{j+1}-e_j)/2}$.
In row $i$ the ratio of successive absolute entries is
$(r_i/\sqrt{r_jr_{j+1}})^{e_{j+1}-e_j}$. It is at least $4$ up to the
diagonal ($j<i$) and at most $1/4$ subsequently ($j\ge i$). The sum of the
off-diagonal magnitudes in a row is at most $2\sum_{m\ge1}4^{-m}=2/3$
times its diagonal magnitude. Strict diagonal dominance implies
invertibility. To see this, take a putative nonzero kernel vector, choose a coordinate of
largest modulus and use the corresponding row. Column scaling is reversible.
For a polynomial with fewer than $t$ terms apply the same result to an
equal-sized subset of the points.
\end{proof}

\begin{theorem}[Uniform support growth]
With the Dirichlet GRH theorem as input, for every sufficiently large
absolute order discriminant $D=|\Delta|$,
\[
                  s(\Delta)\gg\frac{D^{1/4}}{\sqrt{\log D}} .
\]
The implied constant is absolute. Thus each fixed support size occurs at
only finitely many discriminants, including orders of arbitrary conductor.
\end{theorem}
\begin{proof}
Write $D=f^2|d_K|$ for $K=\mathbb Q(\sqrt\Delta)$, with $d_K$ its
fundamental discriminant. Set
$X=\lfloor D^{1/4}\sqrt{\log D}/10\rfloor$ and
$H=\lfloor\log X\rfloor$. Let $S$ consist of the odd primes in $(X/2,X]$
which split in $K$ and do not divide $f$.

\emph{Counting the split primes.}
The quadratic factorization $\zeta_K(s)=\zeta(s)L(s,\chi_{d_K})$
supplies RH for this specific Dedekind zeta function. The uniform prime
ideal theorem of Greni\'e and Molteni, Corollary 1.2, yields, for $y\ge100$,
\[
  |\psi_K(y)-y|
   \ll\sqrt y\big((1+\log y)\log |d_K|+1+\log^2y\big).
\]
Here $\psi_K$ assigns weight $\log N\mathfrak p$ to every power of a prime
ideal $\mathfrak p$. In degree two the contributions other than first powers
of split rational primes are $O(\sqrt y\log y+\log|d_K|)$.
Indeed, over each $p\le\sqrt y$ all prime-power terms have combined
weight $O(\log y)$, and the sum of $\log p$ at first powers of ramified
primes is at most $\log|d_K|$. A split $p$ contributes $2\log p$ at first
powers. Removing the primes dividing $f$ costs another $O(\log(2f))$.
Taking the difference at $y=X$ and $X/2$ therefore gives
$2\sum_{p\in S}\log p=X/2+o(X)$ uniformly in $\Delta$.
In particular $|S|\gg X/\log X$.
See L.~Greni\'e and G.~Molteni,
\href{https://arxiv.org/pdf/1312.4463}{\emph{Explicit versions of the
prime ideal theorem for Dedekind zeta functions under GRH}}.

\emph{Extracting spaced primes.}
For $1\le m\le X$ let $b_m=|S\cap(m,m+H]|$.
Then $\sum b_m=H|S|\gg X$, since each selected prime occurs for exactly
$H$ starts. The unconditional prime second-moment estimate of
Bazzanella, Languasco and Zaccagnini, Lemma 4 with $k=2$, states in our
range that $\sum_{m\le X}(\psi(m+H)-\psi(m))^2\ll X\log^2X$.
As $b_m\log(X/2)\le\psi(m+H)-\psi(m)$, this gives $\sum b_m^2\ll X$.
Cauchy--Schwarz shows that at least a constant times $X$ starts have
$b_m>0$. Take a maximal subset $S'$ of $S$ with pairwise distances at
least $H$. For every occupied start some $p\in S$ lies in $(m,m+H]$
and some $p'\in S'$ has $|p-p'|<H$. Each $p'$ accounts for at most
$3H+1$ such starts. Thus $t=|S'|\gg X/H$.
See D.~Bazzanella, A.~Languasco and A.~Zaccagnini,
\href{https://arxiv.org/pdf/0809.2967}{\emph{Prime numbers in logarithmic
intervals}}, \S3. Only the upper bound for the second moment is used.

\emph{Producing separated conjugates.}
For $p\in S'$, $(\Delta/p)=1$, $p\nmid\Delta$ and $4p^2<D$ for large
$D$. Choose a root of $b^2\equiv\Delta\pmod p$, impose the parity
$b\equiv\Delta\pmod2$, and take the representative $-p<b\le p$.
Then $c=(b^2+D)/(4p)$ is integral, $p<c$, and $[p,b,c]$ is a primitive
reduced form. Its conjugate $z_p$ satisfies
$\big||z_p|-\exp(\pi\sqrt D/p)\big|\le2079$ by the uniform estimate
in Bilu--Luca--Pizarro-Madariaga, \S2. Since $\sqrt D/X\to\infty$,
we have $\log|z_p|=\pi\sqrt D/p+o(1)$ uniformly. For $p<p'$ in
$S'$ the difference of these log-magnitudes is at least
\[
 \frac{\pi\sqrt D\,H}{X^2}-o(1)\longrightarrow 25\pi>\log16
\]
in the sense that the displayed lower bound tends to $25\pi$.
Order the $t$ conjugates by increasing modulus and apply the lemma. Every
annihilating polynomial has at least $t+1\gg X/H\gg D^{1/4}/\sqrt{\log D}$
terms.
\end{proof}

This proof of uniform growth uses two published estimates and an explicit
matrix argument. It does not determine the exceptional discriminants for
four terms. For an elementary variant avoiding the prime second moment,
take $X=D^{1/4}/10$ and use all eligible primes with spacing at least one.
The same proof then yields $s(\Delta)\gg D^{1/4}/\log D$.

There is also an unconditional qualitative comparison. For each fixed
$t$, choose $t$ relatively compact open disks in the modular curve with
positive hyperbolic area, such that every $j$-modulus in each disk is more
than sixteen times every $j$-modulus in the preceding disk. Such disks
exist by continuity near points $iy$ tending to the cusp, where
$j(iy)\sim e^{2\pi y}$. Equidistribution of full CM packets, including
nonmaximal orders, implies that for all sufficiently large $D$ the packet
meets each disk. Indeed one can integrate a nonnegative continuous bump of
positive integral supported in each disk. The lemma then gives
$s(\Delta)>t$. This independently yields qualitative finiteness for bounded
support. The split-prime argument gives the quantitative lower bound
$D^{1/4}/\sqrt{\log D}$. See L.~Clozel and E.~Ullmo,
\href{https://ullmo.ihes.fr/publications/clozel-ullmo.pdf}{\emph{\'Equidistribution
de points de Hecke}}, Theorem 2.4.

\begin{corollary}[A finite certificate with 7,958 conjugates]
For $\Delta=1-2^{61}$, every nonzero rational polynomial vanishing at
$j$ of discriminant $\Delta$ has at least $7{,}959$ terms, independently of
any zero-location theorem.
\end{corollary}
\begin{proof}
Let $D=2^{61}-1$ and $B=\lfloor D^{1/4}\rfloor=38{,}967$.
The accompanying integer sieve counts $7{,}958$ integers $1\le a\le B$
all of whose odd prime factors $p$ have $(\Delta/p)=1$.
For each such $a$, Hensel lifting at odd primes and the Chinese remainder
theorem show that $\Delta$ is a square modulo $4a$. The $2$-part follows
from $\Delta\equiv1\pmod{2^{61}}$. Every $a$ counted is coprime to
$\Delta$. Hence there is a primitive reduced form with leading coefficient
$a$ by Proposition 3.1(1) of Bilu--Luca--Pizarro-Madariaga.
For $a<a'\le B$, the exponents $\pi\sqrt D/a$ and $\pi\sqrt D/a'$
differ by at least $\pi\sqrt D/B^2\ge\pi$. Their exponentials are
both at least $e^{\pi B}$, so the absolute error $2079$ changes their
magnitudes by less than one percent. The actual ratio therefore exceeds
$(0.99/1.01)e^\pi>16$. Apply the separated-size lemma to all $7{,}958$
moduli. The computation needs neither a numerical value of $j$ nor the
coefficients of the class polynomial.
\end{proof}

\subsection{An explicit cutoff for equations with four terms}

\begin{theorem}[Four-term cutoff]
Using Dirichlet GRH as supplied by the upstream result,
$s(\Delta)\ge5$ for every negative order discriminant with
$D=|\Delta|\ge10^{11}$.
\end{theorem}
\begin{proof}
Put $R=\sqrt D$. Seek one odd split prime coprime to the conductor in
each of the three intervals $(uR,vR]$ in the table below. Write
\[
 E(y)=\sqrt y\left[\left(\frac{\log y}{2\pi}+2\right)\log D+
                  2\left(\frac{\log^2 y}{8\pi}+2\right)\right],
 \quad B(y)=2\sqrt y\log y+\log D+2\log2 .
\]
Put $E(0)=0$. Corollary 1.2 of Greni\'e--Molteni gives
$|\psi_K(y)-y|\le E(y)$ for $y\ge100$.
Twice the sum of $\log p$ for the desired primes in $(x,y]$ is at least
$y-x-E(x)-E(y)-B(y)$ when $x=0$ or $x\ge100$.
To see the removal cost, all prime ideal powers over the rational primes
$p\le\sqrt y$ contribute at most $2\sqrt y\log y$. At larger $p$,
inert primes make no contribution. Ramified primes cost at most
$\log|d_K|$, and removing the split primes at the conductor costs at most
$2\log f$. Their sum is $\log D$. The term $2\log 2$ permits the
explicit removal of $2$.

At $D=10^{11}$ the following are outward integer bounds, computed with
$316227<R<316228$ and $\pi>3.14$. The columns correspond to $x=uR$ and
$y=vR$. The accompanying \texttt{verify\_four\_term.py} checks them using
rational arithmetic, including the logarithm bounds.
\[
\begin{array}{c|r|r|r}
 (u,v] & \text{lower bound on }y-x & \text{upper bound on }E(x)+E(y)
      & \text{upper bound on }B(y) \\ \hline
 (0,0.046]       &14546&12138&2340\\
 (0.048,0.189]   &44588&38992&5405\\
 (0.230,0.500]   &85381&75082&9548
\end{array}
\]
All three margins are positive. They persist for larger $R$ after
normalizing by $R$. Indeed $E(cR)/R$ is a positive sum of constants
times $R^{-1/2}$ times products of at most two factors chosen from
$\log R,\log(cR)$. For $c\ge0.046$ these logarithms exceed $9$.
The derivative of the logarithm of each term with respect to $\log R$
is at most $-1/2+2/9<0$. The terms in $B(cR)/R$ decrease by the same
argument, including $R^{-1}\log D$. The normalized interval widths are
constant.

For each of the resulting primes $p_1,p_2,p_3$, Proposition 3.1(1) of
Bilu--Luca--Pizarro-Madariaga gives a reduced form of leading
coefficient $p_i$ since $p_i\nmid D$, $\Delta$ is a square modulo $4p_i$,
and $4p_i^2\le D$. Include the principal form of leading coefficient $1$.
For successive prime intervals ending at $v$ and beginning at $u'$, the
ratio of the associated larger modulus to the smaller one is at least
\[
       \frac{e^{\pi/v}-2079}{e^{\pi/u'}+2079}.
\]
Here $\pi(1/v-1/u')>2.84$ for both gaps and all four exponentials exceed
$e^{13}>207900$. Thus the ratio exceeds $(0.99/1.01)e^{2.84}>16$.
For the principal form versus $p_1\ge3$, the ratio is at least
$(e^{\pi R}-2079)/(e^{\pi R/3}+2079)>16$. The separated-size lemma
applies to the four resulting conjugates. This gives an explicit finite
range for the remaining four-term classification.
\end{proof}

\subsection{Exhausting the intermediate discriminants}
For $10^7\le D<10^{11}$ an integer sieve proves $s(-D)\ge5$ for every order discriminant. Let $a<b$ be two leading coefficients of reduced primitive forms. The following conditions suffice for the absolute values of the corresponding conjugates, in decreasing order, to differ by a factor greater than sixteen:
\[
       D\ge16a^2,\qquad 9D(b-a)^2\ge16a^2b^2.
\]
Indeed $U=\exp(\pi\sqrt D/a)$ then exceeds $34\cdot2079$, and
$U/V\ge e^{4\pi/3}>32$ for $V=\exp(\pi\sqrt D/b)$. Thus
$U-2079>16(V+2079)$. The absolute error bound for $|j|$ proves the claim.

In a bin $L\le D<\min(2L,10^{11})$, consider any odd prime $q$ with $4q^6\le L$. If $-D$ is a nonzero square modulo $q$, the four leading coefficients $1,q,q^2,q^3$ are available. To see this, lift the square root modulo the required odd power, impose the correct parity, and reduce the middle coefficient to absolute value at most the leading coefficient. The constant coefficient is at least the leading coefficient because $4q^6\le D$, and the form is primitive because the middle coefficient is a unit modulo $q$. The consecutive pairs satisfy both integer separation tests. The same four powers of $2$ work whenever $D\equiv7\pmod8$, since an odd number congruent to one modulo eight is a square modulo every power of two. The required size inequalities hold throughout this range.

The checker counts entire residue classes certified by these tests using the Chinese remainder theorem. For the remaining individual discriminants it finds three odd split primes below 1000 which, together with $1$, pass the same tests. Counts are disjoint, and the checker balances each bin against the exact number of $D\equiv0,3\pmod4$ therein.
\[
\begin{array}{l|r}
\toprule
\text{Reason for four separated conjugates}&\text{Discriminants}\\ \midrule
2\text{ splits}&12{,}498{,}750{,}000\\
\text{An odd split prime certifies its first three powers}&37{,}449{,}133{,}804\\
\text{Three separately checked odd split primes}&47{,}116{,}196\\ \midrule
\text{Total}&49{,}995{,}000{,}000\\
\bottomrule
\end{array}
\]
The enumeration and the per-bin counts are in \texttt{checks/cm\_range.cpp} and \texttt{checks/cm\_range.out}. These finite exclusions use no hypothesis about zeros.

\subsection{Exact tests for discriminants through ten million}
\label{sec:cm-small}
Write $D=|\Delta|$. Among the $5{,}000{,}000$ positive
integers $3\le D\le10^7$ with $D\equiv0,3\pmod4$, precisely $67$ have
$h(-D)<4$. For $4{,}999{,}853$ of the $4{,}999{,}933$ others, the following
finite argument proves $s(-D)\ge5$. The $80$ remaining discriminants are
listed with their class numbers in \texttt{checks/cm\_small\_unresolved.txt}.
Their class numbers are $4$ in $50$ cases, $5$ in $29$ cases, and $6$ in the
single case $D=4075$, the largest value left. These are unresolved inputs to
the test. No four-term relation is inferred for any of them.

For four roots with moduli $0<r_0<\cdots<r_3$ and integer exponents
$e_0<\cdots<e_3$, divide each absolute term of the generalized Vandermonde
determinant by the diagonal term $\prod_i r_i^{e_i}$. For a permutation
$\sigma$, the quotient is
\[
 \prod_{k=0}^2
 \left(\frac{\prod_{\sigma(i)>k} r_i}{\prod_{i>k}r_i}\right)^{e_{k+1}-e_k}
 \le \prod_{i=0}^3r_i^{\sigma(i)-i}.
\]
Every bracket is at most one, since the denominator uses the largest radii.
If the $23$ displayed upper bounds for nonidentity permutations sum to less
than one, the determinant is nonzero. This also excludes any relation using
fewer than four powers, since a set of fewer exponents can be enlarged.
When $r_{i+1}\ge4r_i$, the product on the right is at most
$4^{\sum_i i(\sigma(i)-i)}$. Indeed, the exponents of each successive ratio
$r_{k+1}/r_k$ are nonpositive by the same rearrangement. Summing these $23$
rational numbers gives $928541/1048576<1$.

The first pass enumerates primitive reduced forms $[a,b,c]$, taking
$0\le b\le a\le c$ and $D=4ac-b^2$. Negative $b$ gives conjugate moduli
of the same absolute value. If $a<a'$ and
\[
 D\ge16a^2,\qquad 9D(a'-a)^2\ge16a^2a'^2,
\]
then $E_a=\exp(\pi\sqrt D/a)$ satisfies $E_a/E_{a'}>e^4>32$ and
$E_a>e^{12}>34\cdot2079$. Thus $E_a-2079>16(E_{a'}+2079)$.
The uniform cusp estimate for $j$ yields a certified ratio exceeding $16$.
The program initially sieves with $a\le80$ and then enumerates all
$a\le\sqrt{D/3}$ on survivors, stopping once four suitable forms occur.
Class numbers on the remaining inputs use weight one for $b=0$, $b=a$
or $a=c$, and weight two elsewhere. This proves $4{,}999{,}603$ exclusions
and leaves $330$ inputs of class number at least four.

For the second pass we use rigorous enclosures $[L_i,U_i]$ of conjugate
moduli. In addition to the factor-four test, we use every four strictly
ordered disjoint intervals and compare
\[
 \sum_{\sigma\ne\mathrm{id}}\prod_i U_i^{\sigma(i)}
 \quad\hbox{with}\quad \prod_i L_i^i.
\]
A strict upper bound by the latter implies the preceding determinant test.
The factor-four comparison certifies $228$ inputs, and the interval
determinant comparison adds $22$.

All enclosures use integer endpoints with a common denominator $10^{36}$
and outward rounding. For $D\ge100a^2$ the cusp estimate suffices.
Elsewhere put $q=e^{-\pi\sqrt D/a+i\pi b/a}$ and evaluate
$|j|^2=|q|^{-2}|E_4(q)|^6\prod_{n\ge1}|1-q^n|^{-48}$. For reduced forms,
$|q|\le e^{-\pi\sqrt3}<1/200$. Both series and product are truncated
at $K=16$. The error in $E_4$ is at most
$960(K+1)^3/200^{K+1}$, using $\sigma_3(n)<2n^3$ and a geometric tail.
The absolute logarithm of the missing product factor is at most
$192/200^{K+1}$, by $|\log|1-z||\le2|z|$ for $|z|<1/2$.
For such an error $t<1/2$, the omitted factor lies between $1-2t$ and
$1+2t$. Sines and cosines use degree-$50$ Taylor polynomials with error
$4^{51}/51!$. The bounds $1+u\le e^u\le(1-u)^{-1}$ for $0\le u<1$,
applied at scale $2^{-36}$ followed by repeated squaring, enclose the
exponential. The constants enclose $\pi$ between
$3.141592653589793238$ and $3.141592653589793239$. These rational bounds
also follow from Machin's formula and the alternating arctangent series.

\paragraph{The class-number-six case.}
At $\Delta=-4075$ take the three conjugates of forms with $a=25$ and
$b=25,15,5$, in that order, and append the principal conjugate. Call their increasing moduli $r_1,r_2,r_3,r_4$.
The certified enclosures are $2365<r_1<2366$, $2846<r_2<2847$,
$3701<r_3<3703$, $r_4>10^{80}$. Remove the smallest exponent from
a four-term polynomial and let $g_i=e_{i+1}-e_i\ge1$ be the three gaps,
so $e_1=0$. We show that $\det(z_i^{e_j})$ is nonzero for every choice.

For a permutation $\sigma$, the magnitude of its term divided by the
diagonal magnitude equals
\[
 \prod_{j=1}^3(r_j/r_{j+1})^{\sum_{i\le j}(e_{\sigma(i)}-e_i)} .
\]
Each exponent on a radius ratio is a nonnegative linear combination of
the $g_i$, by comparing the selected $j$ exponents with the $j$ smallest.
If $g_1\ge12$ or $g_2\ge12$, the sum over the 23 nonidentity permutations
is bounded above by its value at the appropriate gap vector $(12,1,1)$ or
$(1,12,1)$ with $r_j/r_{j+1}$ replaced by their outward rational bounds.
Both sums are strictly less than one.

It remains to check $1\le g_1,g_2\le11$. Divide all four roots by $4096$,
which rescales the determinant by columns. The contribution from
permutations placing the largest exponent at the principal root, after
dividing by that root's scaled largest power, is the $3\times3$ determinant
of the scaled smaller roots with exponents $(0,g_1,g_1+g_2)$. Its modulus
exceeds $10^{-3}$ in all 121 cases, certified by complex interval
arithmetic. Each of the other 18 permutation terms contains a factor
$r_3/r_4$ in the ratio formula at $j=3$. After the same division the
diagonal magnitude is at most one, so their sum in modulus is below
$18\cdot3703/10^{80}$. This proves the exclusion uniformly in $g_3$.
The certificate \texttt{checks/verify\_4075.py} encloses the complex roots
using the Eisenstein series and discriminant product with explicit tails.

\subsection{The degree-six theorem and the remaining cases}
\begin{theorem}
Every singular modulus of degree at least six requires at least five nonzero coefficients in every nonzero rational polynomial that vanishes at it. For degree at least four, the same conclusion holds except possibly at the following 79 absolute discriminants. The assertion also holds over the corresponding CM field.
\end{theorem}
\begin{center}\small
\begin{tabular}{@{}c p{.84\linewidth}@{}}
\toprule Degree & Remaining absolute discriminants \\ \midrule
4 & 39, 55, 56, 63, 68, 80, 128, 136, 144, 155, 156, 171, 184, 196, 203, 208, 219, 220, 252, 256, 259, 275, 291, 292, 323, 328, 355, 363, 387, 388, 400, 475, 507, 568, 592, 603, 667, 723, 763, 772, 955, 1003, 1027, 1227, 1243, 1387, 1411, 1467, 1507, 1555. \\
5 & 47, 79, 103, 127, 131, 179, 188, 227, 316, 347, 412, 443, 508, 523, 571, 619, 683, 691, 739, 787, 947, 1051, 1123, 1723, 1747, 1867, 2203, 2347, 2683. \\
\bottomrule\end{tabular}\end{center}
\begin{proof}
The analytic cutoff handles $D\ge10^{11}$, and the intermediate sieve covers $10^7\le D<10^{11}$. The small sieve gives precisely the list above together with $4075$ in degree at least four. The preceding complex determinant computation removes that final degree-six case. Linear disjointness gives the assertion over the CM field. The candidate list records where the proofs stop. It does not assert the existence of a four-term relation at any listed discriminant.
\end{proof}

\subsection{How many separated sizes can this argument use?}
There is an upper bound on the strength of the separated-conjugate method itself. For any discriminant $\Delta<-3$ write $D=|\Delta|$. A sequence of conjugates with each absolute value at least sixteen times the preceding one has length $O(D^{1/4})$.
Indeed for a reduced form $[a,b,c]$ with $a\le D^{1/4}$ the conjugate has absolute value $\exp(\pi\sqrt D/a)+O(1)$ with absolute error at most 2079. For sufficiently large $D$ the largest and smallest of these values at any one $a$ have ratio less than two. Thus this part supplies at most $\lfloor D^{1/4}\rfloor$ entries.

For the remaining forms the positive logarithms of the moduli are at most $\pi D^{1/4}+O(1)$. The negative logarithms have absolute value $O(\log D)$. To see the latter bound, a reduced point with $|b|<a$ differs horizontally from either order-three elliptic corner by at least $1/(2a)$. When $|b|=a$ the vertical difference from the closer corner is
\[
 \frac{|D-3a^2|}{2a(\sqrt D+\sqrt3\,a)}\gg D^{-1},
\]
since equality $D=3a^2$ in this case would give the imprimitive multiple of $[1,\pm1,1]$ or discriminant $-3$, and $a\le\sqrt{D/3}$. The triple zero of $j$ at the corners and the positive lower bound off their neighborhoods give $|j|\gg D^{-3}$ (see Bilu--Luca--Pizarro-Madariaga, Lemma 11.6). Packing gaps of length $\log16$ into the resulting interval of logarithms gives the assertion. The power $1/4$ in the lower bound is thus optimal for this fixed-ratio selection method, up to logarithmic factors. This upper bound concerns the proof method. The actual minimum polynomial support is not bounded above by this argument.

\section{The least prime-generating constants for five exponents}
\label{sec:mills}
For $\alpha>1$ call $A>1$ admissible if $\lfloor A^{\alpha^n}\rfloor$ is prime for every integer $n\ge1$. Write $M_\alpha$ for the least admissible $A$ when one exists. Chamberland and Straub give the digits $1.385510850556\ldots$ for $M_{9/4}$ in Example~5.2 of their February 2026 preprint. Their construction produces an admissible number. An earlier prime, discarded by its selection rule, gives a strictly smaller one. Our comparison is with the text of \href{https://arxiv.org/html/2602.22502v1}{arXiv:2602.22502v1, Example 5.2}.

\begin{theorem}\label{thm:mills}
With the ordinary RH supplied by \texttt{riemannHypothesis}, the following are certified initial digits of the least admissible numbers. The last column gives the level of a prime from whose floor interval an infinite prime sequence is proved to continue.
\[
\begin{array}{c@{\qquad}l@{\quad}r}
\toprule
\alpha & M_\alpha \text{ begins with} & \text{level}\\
\midrule
9/4 &1.374277825758088669776655315788864285945977 &7\\
11/5 &1.40371466815791665549209159126764528659 &7\\
21/10 &1.440684278165675686775820360228 &9\\
41/20 &1.478399833482410327000644670 &10\\
201/100 &1.51462541826615615161458 &13\\
\bottomrule
\end{array}
\]
\end{theorem}

\subsection{Where the smaller branch starts}
At level $n$, the floor $p$ means that $A$ belongs to the half-open interval
\[
 I_n(p)=[p^{\alpha^{-n}},(p+1)^{\alpha^{-n}}).
\]
A common way to make $I_{n+1}(p')$ fit inside $I_n(p)$ is to require $p'\ge p^\alpha$ and $p'+1<(p+1)^\alpha$. For nonintegral $\alpha$, the intervals can also overlap when $p'=\lfloor p^\alpha\rfloor<p^\alpha$. An overlap can contain a complete interval at a later level. This is what happens at $\alpha=9/4$.

Any admissible $A$ satisfies $A^\alpha\ge 2$ and $A^{\alpha^2}\ge5$, since $2^\alpha$ is between $4$ and $5$ and $4$ is composite. The floor at the third level can be $37$ since $37<5^\alpha<38$. At the fourth level,
\[
 3455<5^{81/16}<3456<3457<3458<38^{9/4}.
\]
The integers $3455$ and $3456$ are composite and $3457$ is prime. Throughout $3457\le A^{\alpha^4}<3458$ the earlier floors are $2,5,37$. Indeed the lower bound implies $A^{\alpha^2}>5$, the upper gives $A^{\alpha^3}<38$, and $2^\alpha<5$, $38<6^\alpha$, $6<3^\alpha$ give the remaining endpoints. Thus the whole fourth-level interval fits inside the earlier overlap. In particular an infinite continuation from this interval gives $M_{9/4}<38^{\alpha^{-3}}<41^{\alpha^{-3}}$, while the branch in Example~5.2 has third floor $41$. The authors' own all-scales prime-interval theorem at $\alpha=9/4$ already guarantees this continuation. The following proof supplies the much narrower bounds in the table and also works below their uniform exponent threshold.

\subsection{Certifying the lower endpoint and the tail}
Suppose $L$ is a lower bound on every admissible $A$. At level $n$ the next lower bound is
\[
 L' = \max\{L,\,p^{\alpha^{-n}}\},\qquad
 p=\text{least prime at least }\lfloor L^{\alpha^n}\rfloor .
\]
This follows by monotonicity of the floor, with no assumption on the earlier floors of a competitor. For an upper bound, start from the same finite choice of primes and intersect their floor intervals. The certificate checks that the intersection contains a whole terminal interval, that its upper endpoint lies below every earlier upper endpoint, and that this terminal interval admits a continuation. Thus the lower bound and the constructed admissible branch enclose the global infimum. The infimum is itself admissible. Take admissible parameters tending to it from the right. For each fixed $n$, right continuity makes the floors eventually equal to the floor at the limit. Also the initial lower bound $2^{1/\alpha}$ is strictly greater than $1$.

Here is the analytic continuation bound in sufficient detail to reproduce the transition for small $\alpha-2$. Carneiro, Milinovich and Soundararajan prove on RH that for real $y\ge4$ there is a prime in $[y,y+(22/25)\sqrt y\log y]$. For $y=x^\alpha$ and $\alpha>2$, Taylor's theorem gives
$(x+1)^\alpha-1-x^\alpha>\alpha x^{\alpha-1}$ for $x\ge1$. Hence
\[
 x^{(\alpha-2)/2}\ge (22/25)\log x
 \ \Longrightarrow\ 
 \exists\,p\text{ prime},\qquad x^\alpha\le p<(x+1)^\alpha-1
 \quad (x^\alpha\ge4).
\]
This is also Chamberland--Straub's Corollary~3.2. For the five rows take $B=24,30,60,120,800$, respectively. If $h=(\alpha-2)/2$, then $Bh\log 10>1$ and $10^{Bh}>3B$ in each case. The function $x^h/\log x$ increases past $10^B$, and $\log 10<3$, so the implication applies for every $x\ge10^B$. Each prime $p$ past this threshold has a successor whose floor interval lies inside its own, strictly at the upper end. Iteration gives an increasing sequence of lower endpoints with a finite limit. Every old upper endpoint exceeds a later upper endpoint, so this limit lies in every half-open interval. See \href{https://arxiv.org/abs/1708.04122}{Carneiro--Milinovich--Soundararajan, Theorem 5} and \href{https://arxiv.org/html/2602.22502v1}{Chamberland--Straub, Lemma 5.1}.

\subsection{Finite certificates}
In the first seven levels we use the least allowable prime at or above the floor of the current lower bound. Each of the five cases ends this stage with a complete seventh-level interval. If the seventh prime is below $10^B$, we subsequently choose any certified prime with a complete interval inside the previous one until we pass the threshold. These later choices establish an upper bound. Computing the first seven floors of the minimum does not require that every later choice be minimal. Lower and upper bounds for $M_\alpha$ are the two endpoints of the seventh-level interval, with the upper end strict. Its lower endpoint is a global lower bound by construction. A continuation inside it proves the upper bound.

\[
\begin{array}{c|rrrrr|r}
\toprule
\alpha &p_1&p_2&p_3&p_4&p_5& \text{digits of the terminal prime}\\
\midrule
9/4&2&5&37&3457&91637597&41\\
11/5&2&5&37&2819&38923271&37\\
21/10&2&5&29&1213&2993027&126\\
41/20&2&5&29&997&1403869&223\\
201/100&2&5&29&877&823117&1577\\
\bottomrule
\end{array}
\]
The certificate contains every prime, floor and primality witness. If $L=b^{\alpha^{-i}}$ and $\alpha=m/d$, the verifier checks
$f^{d^{n-i}}\le b^{m^{n-i}}<(f+1)^{d^{n-i}}$ to establish a floor $f$. It proves every skipped candidate composite by an integer divisor or a failed strong probable-prime test. Each large prime $P$ comes with a factored divisor $F\mid P-1$, $F^2>P$, and an order witness for each prime factor $q\mid F$. The witness has $a^{P-1}\equiv1\pmod P$ and $\gcd(a^{(P-1)/q}-1,P)=1$. Thus for each prime divisor $r$ of $P$, the order of $a\pmod r$ contains the full $q$-part of $P-1$. It follows that $r\equiv1\pmod F$ for every $r\mid P$. A composite $P$ would have a prime divisor at most $\sqrt P$. The verifier recurses on $q$ and uses trial division below $10^4$. Its decimal enclosures use repeated, outward-rounded rational roots. The table therefore uses neither probabilistic primality nor floating-point comparisons. One branch for each exponent suffices for this argument. Whether every interval between consecutive powers contains a prime at $\alpha=2.01$, or even at $\alpha=2.2$, is a separate question.

\section{A prime plus a multiple of a prime cube or fourth power}
Lee and O'Clarey asked whether every integer greater than 24 is the sum of a prime and a positive multiple of a prime square. Their Theorem~1.3 establishes this under Dirichlet GRH. Campbell also reports a direct unconditional proof in the revision of 24 November 2026. Increasing the exponent to three gives the following different cutoff. \href{https://arxiv.org/html/2605.02426v1}{Lee and O'Clarey}; \href{https://arxiv.org/abs/2609.14714v2}{Campbell}.

\medskip\noindent\textbf{Theorem.} Exactly 62 positive integers fail to have a representation
\[
                  n=p+m r^3,\qquad p,r\text{ prime},\quad m\ge1.
\]
The largest exception is $10\,500$. The complete list is in \texttt{checks/cube.out}.

\smallskip\noindent\textit{Proof.}
Let $q$ be the smallest prime that does not divide $n$, and put $P_q=\prod_{r<q,\ r\ {\rm prime}}r$. Then $P_q\mid n$. If a prime $p<n$ satisfies $p\equiv n\pmod{q^3}$, it supplies the required representation. Corollary~1.2 of Lamzouri, Li and Soundararajan supplies such a congruent prime with
\[
                       p\le [\varphi(q^3)\log(q^3)]^2
\]
under Dirichlet GRH. This is the precise input supplied by the Lean theorem, together with the published implication. \href{https://arxiv.org/abs/1309.3595}{Lamzouri, Li and Soundararajan}.

For $q\ge43$ the displayed upper bound is less than $q^8$ since $3\log q<q$. Moreover, $P_q>q^8$ for all such primes. The checker compares integers for $43\le q\le257$. Thereafter Bertrand's postulate puts the next prime $s$ below $2q$, so $P_s=qP_q>q\cdot q^8>(2q)^8>s^8$. Consequently $p<n$ throughout this range.

For $q\le41$ let $M_q$ be the maximum of the least prime in each reduced class modulo $q^3$. A sieve through eight million supplies a representative of every class for each of these moduli. As soon as all $q^2(q-1)$ classes appear, the last prime used equals $M_q$. Every $n>M_q$ with this least nondivisor $q$ is covered. The remaining candidates are the positive multiples of $P_q$ up through $M_q$ that are coprime to $q$. For each candidate the checker visits every prime cube below $n$ until it finds a positive multiple with prime complement or exhausts the possibilities.
\[
\begin{array}{c|rrrrrrrr}
q&2&3&5&7&11&13&17&19\\ \hline
M_q&17&109&1621&9371&47737&120401&349529&540961\\
\text{candidates}&9&36&216&268&207&48&11&1\\
\text{exceptions}&6&11&25&15&5&0&0&0
\end{array}
\]
For $q=23,29,31,37,41$ there are no remaining candidates. This accounts for every positive integer. The finite calculation uses only integer arithmetic.

The last exception can also be checked by hand. For $n=10\,500$, a cube with base $r=2,3,5,7$ forces $p=r$, and $r^3\nmid n-r$ in each case. The other possible bases are $11,13,17,19$. In order of increasing positive multiplier $m$, proper divisors of all positive complements $n-mr^3$ are, respectively,
\[
(53,2,3,2,5,2,7),\qquad(19,2,3,2),\qquad(37,2),\qquad(11).
\]
This proves sharpness.\hfill$\square$
\subsection*{Fourth powers}
The same argument gives an exact cutoff when $r^3$ is replaced by $r^4$.
\textbf{Exactly 811 positive integers fail to be a prime plus a positive multiple of a prime fourth power. The largest is $1\,561\,560$.}

Retain $q$ and $P_q$ from the preceding proof. For $q\ge53$, the least congruent prime modulo $q^4$ is at most $[4\varphi(q^4)\log q]^2<q^{10}$. Exact comparisons show $P_q>q^{10}$ for all primes from 53 through 1031. Bertrand's postulate propagates the inequality from any $q>2^{10}$ to its successor. Thus $n$ exceeds the least congruent prime for the entire tail.

For $q\le23$, enumerate all reduced classes modulo $q^4$ as above, then test the remaining multiples of $P_q$. The corresponding maxima of the least primes, in increasing order of $q$, are
\[
 41,\ 647,\ 20063,\ 92369,\ 1189213,\ 2987429,\ 9951083,\ 19423919,\ 42026629.
\]
The nine groups yield, respectively, $12,48,230,255,212,46,7,1,0$ exceptions. When testing another prime base $r$ that divides $n$, the only possible prime complement is $r$ itself. This observation shortens the finite check.

For the six remaining primes $29\le q\le47$, we have $4\log q<16$. Hence the least-prime theorem covers $n>B_q=[16q^3(q-1)]^2$. For each positive multiple $n$ of $P_q$ through $B_q$ with $q\nmid n$, the checker exhibits a prime in the residue class of $n$ modulo $q^4$ strictly below $n$. The numbers of inputs in these six groups are $516676,30586,4130,256,9,0$. Sieve lookup and trial division certify all the primes. The full exceptional list and reproducible computation are in \texttt{checks/fourth.out} and \texttt{checks/fourth.py}.

As an independent finite check, for every prime $r$ with $r^4\le2\,000\,000$, one can take the first prime $p$ in every occupied class modulo $r^4$ and mark $p+r^4,p+2r^4,\ldots$ through two million. We obtained the identical list of 811 unmarked positive integers by this procedure. The primorial argument above certifies completeness beyond this finite check.

\section{A one-sided test on sparse twists and sparse scales}
\label{sec:sparse}
A separate consequence uses the proof's zero-detection argument before its
final GRH theorem. The published Lean criterion takes an eventual inequality
for every auxiliary prime in one residue class and every sufficiently large
real scale. Both sets of inputs can be thinned. The conclusion below is a
reduction proved here on paper. Sandi Setiawan already gives the fourth-power
mesh and a chord interpolation argument for the ordinary zeta prime sum
(\emph{The Primacohedron Trace Formulation of the Riemann Hypothesis},
22 November 2026, Theorems 3.2 and 4.1,
\href{https://www.academia.edu/181628231}{author's manuscript}).
Here the additional reduction selects a thin family of quadratic twists
using the repository's quantitative zero-density bound. We make no claim
of priority for the combined criterion.

Fix an integer $Q>0$ divisible by $8$ and a reduced residue class
$a\pmod Q$ with $a\equiv7\pmod8$. For primes $d\equiv a\pmod Q$, put
\[
 w_d(n)=\mathbf1_{n\equiv1\bmod Q}\,\Lambda(n)
       \bigl(1-(n/d)\bigr),\qquad
 R_d(x)=\sum_{1\le n\le x} w_d(n)(1/n-1/x)
       -\frac{\log x}{\varphi(Q)}-c_{d,Q}.
\]
Here $(n/d)$ is the Legendre symbol and $c_{d,Q}$ is the unique constant
for which $R_d(x)\to0$. The repository gives its explicit logarithmic
derivative formula in \texttt{New.logDerivMean} and identifies it in
\texttt{New.primeRemainder}. All $w_d(n)$ are nonnegative.

\begin{theorem}[Sparse one-sided zero detection]
Fix $1/2<\beta<1$. Let $S$ be a set of the above primes and
$S(X)=\#\{d\in S:d\le X\}$. Suppose
\[
 \limsup_{X\to\infty}\frac{\log(1+S(X))}{\log X}
 \ \ge\ 1-\frac{\beta-1/2}{18}.
 \tag{S}
\]
For each $d\in S$, suppose there is an increasing unbounded sequence
$x_j\ge1$, allowed to depend on $d$, with
\[
 x_{j+1}-x_j=O_d(x_j^{(1+\beta)/2}),\qquad
 R_d(x_j)\le C_d x_j^{\beta-1}\quad(j\text{ sufficiently large}),
 \tag{E}
\]
where $C_d>0$. For every Dirichlet character $\psi\bmod Q$,
$L(s,\psi)$ is nonzero in $\Re s>\beta$, with the principal pole excepted.
\end{theorem}

\begin{proof}
First recover all intervening scales. Write $t=1/x$. On any compact interval
$[u,v]\subset(0,\infty)$ the function
\[
 F(t)=\sum_{n\ge1}w_d(n)\,(1/n-t)_+
\]
is a finite sum of convex functions, hence lies below its chord. We have
$R_d(1/t)=F(t)+(\log t)/\varphi(Q)-c_{d,Q}$. Write $I f$ for the affine
interpolant from $u,v$. Since $t\mapsto\log t+t^2/(2u^2)$ is convex on
$[u,v]$,
\[
 0\le\log t-I(\log)(t)
 \le\frac{(t-u)(v-t)}{2u^2}
 \le\frac{(v-u)^2}{8u^2}.
\]
Take $u=1/x_{j+1}, v=1/x_j$. The ratio $(v-u)/u$ equals
$(x_{j+1}-x_j)/x_j$, so the interpolation loss is
$O_d(x_j^{\beta-1})$. Moreover $x_{j+1}/x_j\to1$. Applying (E) at
the endpoints gives $R_d(x)\le K_d x^{\beta-1}$ throughout the tail for
some $K_d>0$.

Suppose $L(\rho,\psi)=0$ with $\sigma=\Re\rho>\beta$. At $\rho=1$
the classical Dirichlet nonvanishing theorem applies to nonprincipal
characters. The ordinary zero-free region $\Re s\ge1$ therefore leaves
$\sigma<1$. Choose $\gamma$ with $\beta<\gamma<\sigma$ and
set $H=|\Im\rho|$. Call an admissible prime $d$ bad if at least one
$L(s,\xi(\cdot/d))$, $\xi\bmod Q$, has a zero in
$\Re s\ge\gamma$, $|\Im s|\le H$. The upstream lemma
\texttt{New.zero\_density\_quadTwist}, with
$\delta=(\gamma-1/2)/9$, gives
\[
 \#\{d\le X:d\text{ bad}\}
   = O_{Q,a,\gamma,H}(X^{1-(\gamma-1/2)/18}).
\]
This exponent is strictly smaller than the right side of (S).
Thus some $d\in S$ lies outside the bad set. Fix this $d$.
The real-variable estimate already proved implies
$R_d(x)\le x^{\gamma-1}$ eventually, since $K_d x^{\beta-\gamma}\to0$.

For clarity, here is where one good twist suffices. The meromorphic
transform of $v\mapsto R_d(e^v)$ in \texttt{New.isPolar\_primeRemainder}
has residue at $s\ne1$ equal to
\[
 \frac{m(s)-m_d(s)}{\varphi(Q)s(1-s)},
\]
where $m$ and $m_d$ sum the zero multiplicities of the pole-removed
untwisted and twisted functions over all $\xi\bmod Q$. The transform is
regular at $1$. In the box, $m_d=0$. On its real segment
$\gamma<s<1$ the residues are nonnegative; for real $s\ge1$ they
vanish. The exponential upper bound $R_d(e^v)\le e^{(\gamma-1)v}$
and these real signs are precisely the hypotheses of the repository's
one-sided Landau lemma, \texttt{IsPolar.landau}. It forces the residue
at $\rho$ to vanish. Then $m(\rho)=0$, contrary to the stipulated zero of
$L(s,\psi)$. This is the fixed-$d$ part of the proof of
\texttt{New.LFunction\_ne\_zero\_of\_primeRemainder\_le}.
\end{proof}

\paragraph{What a counterexample would force.}
If a zero of real part $\sigma>\beta$ existed, let $T_\beta(X)$ count
admissible primes $d\le X$ for which
$\limsup_{j\to\infty}j^{4(1-\beta)}R_d(j^4)<+\infty$.
As $j^4$ meets (E)'s gap restriction, the proof actually gives
\[
 T_\beta(X)=O_{Q,a,\rho,\gamma}
 \bigl(X^{1-(\gamma-1/2)/18}\bigr)
 \qquad(\beta<\gamma<\sigma).
\]
Indeed, every such $d$ would have to be bad in that single box.
Thus the discrepancy at fourth powers would have arbitrarily large
positive normalized excursions for every admitted prime outside
an exceptional set of power-saving size.

\paragraph{An explicit sampling schedule.}
Enumerate the primes in the prescribed class as $p_1<p_2<\cdots$.
Choose $d_k=p_{\lfloor k\log^2(k+2)\rfloor}$ for all sufficiently large
$k$. The prime number theorem in progressions gives
$d_k\sim\varphi(Q)k\log^3 k$. Their reciprocals have a convergent sum,
and their counting exponent in (S) is $1$. Take $x_j=j^4$.
The gap is $O(x_j^{3/4})$. If for every selected $d_k$ and every
$\varepsilon>0$ the sampled values satisfy
\[
 R_{d_k}(j^4)\le C_{k,\varepsilon}\,j^{-2+4\varepsilon}
 \quad\text{for all sufficiently large }j,
\]
apply the theorem with $\beta=1/2+\varepsilon$ for arbitrarily small
positive $\varepsilon$. This excludes all zeros to the right of $1/2$
modulo $Q$ apart from the principal pole. For arbitrary original modulus
$q$, take $Q=18480q$, $a=Q/2-1$; this class meets the conditions because
$Q$ is divisible by $16$ and $\gcd(Q/2-1,Q)=1$. Changing the level adds
finite Euler factors, so the nonvanishing descends to $q$. For these
parameters the upstream one-sided estimate already implies the sampled
inequalities. Their interest is the weaker input an alternative proof
would have to establish. Bounds for finitely many twists or finitely many
scales do not give the required asymptotic estimates.


\section{Further research questions}
\begin{enumerate}
\item \textbf{Four-term equations.} Which discriminants with $h(\Delta)\ge4$ admit a singular modulus satisfying a rational four-term equation? Only 50 discriminants of degree four and 29 of degree five remain, with $|\Delta|\le2683$. Can the complex or nonarchimedean phases distinguish the powers when complex-conjugate roots have equal absolute value? A possible additional test uses the Newton polygon of the class polynomial. Four distinct root valuations at the same prime require at least four slopes in any annihilator and hence at least five nonzero coefficients. The positive-real conjugate test already excludes every $-4n$ with $\omega(n)\ge3$.
\item \textbf{The size of $s(\Delta)$.} Can one recover the remaining factor $\sqrt{\log D}$ using composite leading coefficients in reduced forms? To exceed the $1/4$ power using absolute sizes alone requires an argument capable of exploiting conjugates at comparable radii. Their complex arguments offer additional information.
\item \textbf{Prime-generating constants near two.} Certificates here treat selected rational exponents. Can a common finite construction control a whole interval of exponents, including the changes where a powered endpoint crosses an integer? At exactly two the factor $x^{(\alpha-2)/2}$ no longer grows. The method here does not decide whether the least-prime-above-the-square recursion runs forever within the required intervals.
\item \textbf{Additive cutoffs.} For each fixed $k\ge2$, the least prime not dividing $n$ reduces $n=p+mr^k$ to a prime in a reduced class modulo $q^k$. Its GRH least-prime bound is below $q^{2k+2}$ eventually, and the primorial outgrows this power. The finite residue-class search can thus be repeated for larger $k$. What controls the growth of the largest exception with $k$?
\item \textbf{A nonabelian analytic input.} To transfer the sparse one-sided test to further $L$-functions one needs compatible nonnegative weights, control of the residues, and a zero-density estimate in the corresponding twisting family. Even for the splitting field of $X^3-2$, Dirichlet GRH controls only its abelian characters directly. A bound for the two-dimensional factor would address a genuinely additional zero problem.
\end{enumerate}
\end{document}
